%tra le quadre le dimensioni del carattere (di base = 10)
\documentclass[12]{article}
\usepackage{graphicx} % Required for inserting images
\usepackage{hyperref} %per gli ipertesti 
\usepackage{lineno} %per il numero di linee
\usepackage{comment} %per mettere commenti 
\usepackage{natbib}

\title{Il mio primo documento in LATEX}
\author{Nicola Arcangeli}

%\date{}

%se si mette un begin serve anche un end
\begin{document}

%per settare il titolo
\maketitle

%sommario con ipertesti automatici
\tableofcontents

%lista delle figure
\listoffigures

%inserisce una sezione
\section{Introduzione} \label{sec:intro}
In questo elaborato vengono presentati i risultati di uno studio fitosociologico sulle 
tipologie forestali presenti sul territorio altoatesino. L’obiettivo primo è quello di individuare, 
quantificare e confrontare la diversità delle foreste e procedere con le relative analisi della 
%aggiunge un'interlinea
\smallskip %possibile anche \bigskip
diversità vegetale e funzionale degli hotspot. Per l’analisi della vegetazione e in particolar modo 
per la tipizzazione delle foreste dell’Alto Adige ci si è avvalsi di un database comprendente 924 
rilievi e 1241 specie di piante. Nella scelta delle tipologie forestali sono stati esaminati tre tipi 
di vegetazione: Lariceti, Larici-Cembreti e Cembreti puri.

%sta roba non va più bene ma nel dubbio meglio tenerla

\begin{comment}
I canali di bonifica si sono rivelati essere un luogo fondamentale pe
la conservazione delle specie e degli ecosistemi tipici delle zone umide interne (Bolpagni et al. 2019, Cantonati et al. 
2020, van Rees et al. 2020) che sono tra i più minacciati al mondo (Reid et al. 2019, Dudgeon 
2019). Inoltre la rete dei canali di bonifica è essenziale perché mette in relazione le zone umide 
piccole e isolate, che sono quelle a maggiore rischio di estinzione (Deane et al. 2017). I risultati 
dello studio di Deane et al. (2017) hanno rafforzato questa idea dimostrando che, a parità di 
perdita di superficie, la scomparsa di intere zone umide, anche se piccole, comporta un rischio 
di estinzione per le specie molto più elevato rispetto alla perdita della stessa superficie suddivisa
\end{comment}




\section{Metodi}
In questa tesi è stata utilizzata l'equazione di Newton, raccolta nell'equazione \ref{eq:newton} %automatizza il riferimento

%scriviamo una formula, \frac per far vedere la frazione
% underscore(_) per i pedici
%sqrt per la radice quadrata, tra [] la potenza della radice
\begin{equation}
F=G
\end{equation}

\begin{equation}
F=G\times\frac{\sqrt{m_1\times \sqrt[3]{m_2}}}{r^2}
%mettiamo una label
\label{eq:newton}
\end{equation}

Come visto nella sezione \ref{sec:intro} %per linkare direttamente l'introduzione

Per la tesi sono stati usati
%elenchi puntati
\begin{itemize}
    \item analisi multitemporale
    \item analisi multivariata
    \item analisi multispazzolata
\noindent
\end{itemize}
%elenco numerato
\begin{enumerate}
   \item analisi multitemporale
    \item analisi multivariata
    \item analisi multispazzolata
\end{enumerate}


\section{Risultati}

Sulla base delle metodologie utilizzate abbiamo ottenuto i risultati riportati in Figura \ref{fig:zone} %importante fare le ref per non avere problemi in futuro

%inserisco figura
\begin{figure}
    \centering
   %funzione che inserisce la figura
    \includegraphics[width=\linewidth]{Zoniz. PAAM.png}
    \caption{Bella Zonizzazione}
    \label{fig:zone}
\end{figure}

Ma anche la figura \ref{fig:mappa} non è male eh

\begin{figure}
    \centering
    \includegraphics[width=0.5\linewidth]{Mappona.png}
    \caption{Mappona}
    \label{fig:mappa}
\end{figure}

\section{Discussione}

Protected areas (PAs) often use zoning schemes to balance biodiversity conservation and sustainable human 
use, typically dividing the PA into fully protected zones and others permitting varying \citep{Maeda2005}levels of activity, in line 
with UNESCO and International Union for Conservation of Nature visions. However, the effectiveness of 
zoning in protecting biodiversity is rarely assessed. We reviewed the scientific literature, screening 3,783 ar
ticles and identifying only 60 that specifically evaluated zoning effects on biodiversity. We observed \citep{Miller2017}
geographic, environmental, and taxonomic biases: there are few studies in regions where zoned PAs are 
widely implemented (e.g., Europe); 75% focused on marine environments, none in freshwater, and most tar
geted vertebrates, particularly fish. Biodiversity was mainly measured through species diversity and abun
dance, often from a fisheries management perspective. Critically, few studies assessed biodiversity adapt
ability to global change—a key 2050 Global Biodiversity Framework target. This lack of comprehensive 
assessment hinders informed conservation decisions. We propose a blueprint for multifaceted biodiversity 
monitoring to support holistic evaluations of zoning effectiveness
%test per la bibiliografia
(Rocchini et al., 2025) %non funzionale
\citep{Uyen2025}


\begin{thebibliography}{999}
%[tra le quadre come deve comparire] {tra le graffe la "label"}
% usare \& al posto di &
\bibitem[Miller, J. D. (2017)]{Miller2017}
Miller, J. D. (2017). Reproduction in sea turtles. The Biology of Sea Turtles, Volume I, 51-81.
%
\bibitem[Uyen et al. (2025)]{Uyen2025} 
Uyen, B. T., Marchetto, E., Tordoni, E., \& Rocchini, D. (2025). From virtuality-to-reality: testing the effects of occurrence-based temporal settings on species distribution modeling prediction. Community Ecology, 1-12.

\bibitem[Maeda et al. (2005)]{Maeda2005}
Maeda, T., Hobbs, R. M., \& Pandolfi, P. P. (2005). The transcription factor Pokemon: a new key player in cancer pathogenesis. Cancer research, 65(19), 8575-8578.
\end{thebibliography}


\end{document}
